\section{Discussion and Limitations}
\label{sec:discussion}

\textbf{Implications.} The GF-Score reveals that aggregate robustness metrics systematically hide class-level disparities that matter for deployment. The consistent vulnerability of ``cat'' across 76\% of CIFAR-10 models, and the existence of ImageNet models with $\mathrm{WCR} = 0$ (zero certified robustness on at least one class), demonstrate that passing an aggregate robustness audit provides no guarantee for individual classes. The positive correlation between aggregate robustness and RDI suggests that improving overall robustness through current training methods may inadvertently worsen fairness, reinforcing the tension identified by \citet{xu2021robustfairfairnessadversarial} and \citet{benz2021robustnessoddsfairnessempirical} with new certified evidence. Our framework enables practitioners to detect such issues post-hoc, without retraining, making it complementary to training-time fairness interventions~\citep{wei2023cfaclasswisecalibratedfair, zhang2024fairnessawareadversariallearning}.

\textbf{Limitations.} Our framework inherits the assumptions of the GREAT Score: (1) the generative model must approximate the true data distribution well enough for the certified bound to be meaningful; (2) the local score's certification relies on the confidence margin, which may be loose for models with poorly calibrated softmax outputs. The self-calibration procedure assumes a monotonic relationship between robustness and clean accuracy, which holds within a defense family but may not hold across fundamentally different architectures. Our ImageNet evaluation covers only 5 models due to computational constraints on logit extraction for 50K images across 1000 classes; scaling to more models would strengthen the findings. Finally, while we provide concentration bounds for RDI, tighter bounds for NRGC and WCR remain open theoretical questions.

\textbf{Conclusion.} We introduced the GF-Score, a framework that decomposes the certified GREAT Score into per-class robustness profiles and quantifies their disparity through four metrics grounded in welfare economics. The decomposition is provably exact and comes with finite-sample concentration guarantees. Our attack-free self-calibration eliminates the need for adversarial attacks entirely, achieving rank correlations of $\rho = 0.871$ on CIFAR-10 and $\rho = 1.000$ on ImageNet with RobustBench. Evaluating 22 models across two benchmarks, we find that class vulnerability is remarkably consistent (``cat'' is worst in 76\% of CIFAR-10 models) and that more robust models exhibit greater class-level disparity. These findings demonstrate that aggregate robustness scores are insufficient for safety-critical deployment and that post-hoc fairness auditing of robustness is both feasible and necessary. We release an interactive auditing dashboard alongside our evaluation code to support adoption by practitioners and researchers.
